% GENERATED FILE. Edit canonical lessons, then run npm run generate:books.
% canonical-source-hash: fnv1a64:2bace2f5de38d8e6
% canonical-lessons: PA-C306-chhajj, PA-C306-chhattna, PA-C306-turi, PA-C306-parali, PA-C306-bori

\chapter{Grain and Straw}
\label{ch:pa-a2-grain-and-straw}
\hlchaptermodality{306}

\begin{chapteropening}
\textbf{By the end of this chapter:} \emph{I can talk about the grain and the straw.}

Complete the last lesson of chapter 306: I can talk about the grain and the straw.
\end{chapteropening}

\section[\texorpdfstring{\pa{ਛੱਜ} (chhajj)}{ਛੱਜ (chhajj)}]{\pa{ਛੱਜ} (chhajj) — a winnowing basket}

\label{lesson:PA-C306-chhajj}



\begin{quote}
Before the new one: say the Punjabi for to thresh, then the Punjabi for a heap.
\end{quote}

\subsection*{You'll want to know: \pa{ਛੱਜ}}
\textbf{\pa{ਛੱਜ}} — \emph{chhajj} — \textquotedblleft{}a winnowing basket\textquotedblright{}.

\begin{quote}
\pa{ਮਾਤਾ} \pa{ਜੀ} \pa{ਛੱਜ} \pa{ਨਾਲ} \pa{ਦਾਣੇ} \pa{ਸਾਫ਼} \pa{ਕਰਦੇ} \pa{ਹਨ।}
\end{quote}

\begin{quote}
\emph{mātā jī chhajj nāl dāṇe sāf karde han.}
\end{quote}

\begin{quote}
Mother cleans the grain with a winnowing basket.
\end{quote}

A masculine noun.

\begin{cousinweb}
\textbf{\pa{ਛੱਜ}} (\emph{chhajj}) is inherited from a reconstructed \emph{*chajja}, \textquotedblleft{}basket\textquotedblright{}; CDIAL lists Punjabi \emph{chajj}, \textquotedblleft{}winnowing basket\textquotedblright{}, under it (CDIAL 4964).
\end{cousinweb}

\begin{grammarlens}[title={What you've built}]
One more word for this chapter.
\end{grammarlens}

\subsection*{Your turn}
\begin{itemize}
\raggedright
  \item \emph{Say it:} \emph{chhajj}
  \item \emph{Say it:} \emph{chhajj}, once more
  \item \emph{Say it:} \emph{mātā jī chhajj nāl dāṇe sāf karde han.} — Mother cleans the grain with a winnowing basket.
  \item \emph{Recall:} say the Punjabi for to thresh, then the Punjabi for a heap, then say \emph{chhajj} again
\end{itemize}

\begin{tcolorbox}[breakable,colback=teal!4,colframe=teal!35!black,title={Before you move on}]
What does \pa{ਛੱਜ} mean? (\textquotedblleft{}A winnowing basket\textquotedblright{}.) Say it once more.
\end{tcolorbox}
\section[\texorpdfstring{\pa{ਛੱਟਣਾ} (chhaṭṭṇā)}{ਛੱਟਣਾ (chhaṭṭṇā)}]{\pa{ਛੱਟਣਾ} (chhaṭṭṇā) — to winnow}

\label{lesson:PA-C306-chhattna}



\begin{quote}
Before the new one: say the Punjabi for a heap, then the Punjabi for a winnowing basket.
\end{quote}

\subsection*{You'll want to know: \pa{ਛੱਟਣਾ}}
\textbf{\pa{ਛੱਟਣਾ}} — \emph{chhaṭṭṇā} — \textquotedblleft{}to winnow\textquotedblright{}.

\begin{quote}
\pa{ਅਸੀਂ} \pa{ਹਵਾ} \pa{ਵਿੱਚ} \pa{ਦਾਣੇ} \pa{ਛੱਟਾਂਗੇ।}
\end{quote}

\begin{quote}
\emph{asīṁ havā vich dāṇe chhaṭṭāṅge.}
\end{quote}

\begin{quote}
We will winnow the grain in the wind.
\end{quote}

Past (I …): \pa{ਮੈਂ} \pa{ਦਾਣੇ} \pa{ਛੱਟੇ}, \emph{maiṁ dāṇe chhaṭṭe}.

\begin{cousinweb}
\textbf{\pa{ਛੱਟਣਾ}} (\emph{chhaṭṭṇā}) is inherited from a reconstructed \emph{*chaṭṭ}, \textquotedblleft{}scatter\textquotedblright{}; CDIAL lists Punjabi \emph{chaṭṭṇā}, \textquotedblleft{}to winnow\textquotedblright{}, under it (CDIAL 4968).
\end{cousinweb}

\begin{grammarlens}[title={What you've built}]
One more word for this chapter.
\end{grammarlens}

\subsection*{Your turn}
\begin{itemize}
\raggedright
  \item \emph{Say it:} \emph{chhaṭṭṇā}
  \item \emph{Say it:} \emph{chhaṭṭṇā}, once more
  \item \emph{Say it:} \emph{asīṁ havā vich dāṇe chhaṭṭāṅge.} — We will winnow the grain in the wind.
  \item \emph{Recall:} say the Punjabi for a heap, then the Punjabi for a winnowing basket, then say \emph{chhaṭṭṇā} again
  \item \emph{Write it:} copy \textbf{\pa{ਛੱਟਣਾ}} once, checking the whole word
\end{itemize}

\begin{tcolorbox}[breakable,colback=teal!4,colframe=teal!35!black,title={Before you move on}]
What does \pa{ਛੱਟਣਾ} mean? (\textquotedblleft{}To winnow\textquotedblright{}.) Say it once more.
\end{tcolorbox}
\section[\texorpdfstring{\pa{ਤੂੜੀ} (tūṛī)}{ਤੂੜੀ (tūṛī)}]{\pa{ਤੂੜੀ} (tūṛī) — chopped straw; chaff}

\label{lesson:PA-C306-turi}



\begin{quote}
Before the new one: say the Punjabi for a winnowing basket, then the Punjabi for to winnow.
\end{quote}

\subsection*{You'll want to know: \pa{ਤੂੜੀ}}
\textbf{\pa{ਤੂੜੀ}} — \emph{tūṛī} — \textquotedblleft{}chopped straw; chaff\textquotedblright{}.

\begin{quote}
\pa{ਗਾਂ} \pa{ਤੂੜੀ} \pa{ਖਾਂਦੀ} \pa{ਹੈ।}
\end{quote}

\begin{quote}
\emph{gāṁ tūṛī khāndī hai.}
\end{quote}

\begin{quote}
The cow eats chopped straw.
\end{quote}

A feminine noun.

\begin{cousinweb}
\textbf{\pa{ਤੂੜੀ}} (\emph{tūṛī}) goes back to Sanskrit \emph{tuṣa}, \textquotedblleft{}chaff of grain\textquotedblright{}; CDIAL lists Punjabi only \emph{tuh}, \textquotedblleft{}chaff\textquotedblright{}, under it; the form \emph{tūṛī}, \textquotedblleft{}straw, chaff\textquotedblright{}, is the one it gives for Lahnda (CDIAL 5892).
\end{cousinweb}

\begin{grammarlens}[title={What you've built}]
One more word for this chapter.
\end{grammarlens}

\subsection*{Your turn}
\begin{itemize}
\raggedright
  \item \emph{Say it:} \emph{tūṛī}
  \item \emph{Say it:} \emph{tūṛī}, once more
  \item \emph{Say it:} \emph{gāṁ tūṛī khāndī hai.} — The cow eats chopped straw.
  \item \emph{Recall:} say the Punjabi for a winnowing basket, then the Punjabi for to winnow, then say \emph{tūṛī} again
  \item \emph{Write it:} copy \textbf{\pa{ਤੂੜੀ}} once, checking the whole word
\end{itemize}

\begin{tcolorbox}[breakable,colback=teal!4,colframe=teal!35!black,title={Before you move on}]
What does \pa{ਤੂੜੀ} mean? (\textquotedblleft{}Chopped straw; chaff\textquotedblright{}.) Say it once more.
\end{tcolorbox}
\section[\texorpdfstring{\pa{ਪਰਾਲੀ} (parālī)}{ਪਰਾਲੀ (parālī)}]{\pa{ਪਰਾਲੀ} (parālī) — rice straw}

\label{lesson:PA-C306-parali}



\begin{quote}
Before the new one: say the Punjabi for to winnow, then the Punjabi for chopped straw.
\end{quote}

\subsection*{You'll want to know: \pa{ਪਰਾਲੀ}}
\textbf{\pa{ਪਰਾਲੀ}} — \emph{parālī} — \textquotedblleft{}rice straw\textquotedblright{}.

\begin{quote}
\pa{ਸਿਆਲ} \pa{ਵਿੱਚ} \pa{ਅਸੀਂ} \pa{ਪਰਾਲੀ} \pa{ਉੱਤੇ} \pa{ਸੌਂਦੇ} \pa{ਹਾਂ।}
\end{quote}

\begin{quote}
\emph{siāl vich asīṁ parālī utte saunde hāṁ.}
\end{quote}

\begin{quote}
In winter we sleep on rice straw.
\end{quote}

A feminine noun.

\begin{cousinweb}
\textbf{\pa{ਪਰਾਲੀ}} (\emph{parālī}) is inherited from Sanskrit \emph{palāla}, \textquotedblleft{}stalk, straw\textquotedblright{}; CDIAL lists Punjabi \emph{parālī}, \textquotedblleft{}straw of rice or millet\textquotedblright{}, under it (CDIAL 7958).
\end{cousinweb}

\begin{grammarlens}[title={What you've built}]
One more word for this chapter.
\end{grammarlens}

\subsection*{Your turn}
\begin{itemize}
\raggedright
  \item \emph{Say it:} \emph{parālī}
  \item \emph{Say it:} \emph{parālī}, once more
  \item \emph{Say it:} \emph{siāl vich asīṁ parālī utte saunde hāṁ.} — In winter we sleep on rice straw.
  \item \emph{Recall:} say the Punjabi for to winnow, then the Punjabi for chopped straw, then say \emph{parālī} again
  \item \emph{Write it:} copy \textbf{\pa{ਪਰਾਲੀ}} once, checking the whole word
\end{itemize}

\begin{tcolorbox}[breakable,colback=teal!4,colframe=teal!35!black,title={Before you move on}]
What does \pa{ਪਰਾਲੀ} mean? (\textquotedblleft{}Rice straw\textquotedblright{}.) Say it once more.
\end{tcolorbox}
\section[\texorpdfstring{\pa{ਬੋਰੀ} (borī)}{ਬੋਰੀ (borī)}]{\pa{ਬੋਰੀ} (borī) — a sack}

\label{lesson:PA-C306-bori}



\begin{quote}
Before the new one: say the Punjabi for chopped straw, then the Punjabi for rice straw.
\end{quote}

\subsection*{You'll want to know: \pa{ਬੋਰੀ}}
\textbf{\pa{ਬੋਰੀ}} — \emph{borī} — \textquotedblleft{}a sack\textquotedblright{}.

\begin{quote}
\pa{ਇੱਕ} \pa{ਬੋਰੀ} \pa{ਵਿੱਚ} \pa{ਬਹੁਤ} \pa{ਕਣਕ} \pa{ਹੁੰਦੀ} \pa{ਹੈ।}
\end{quote}

\begin{quote}
\emph{ikk borī vich bahut kaṇak hundī hai.}
\end{quote}

\begin{quote}
A sack holds a lot of wheat.
\end{quote}

A feminine noun.

\begin{cousinweb}
\textbf{\pa{ਬੋਰੀ}} (\emph{borī}) is inherited from a reconstructed \emph{*bōra}, \textquotedblleft{}sacking, sack\textquotedblright{}; CDIAL lists Punjabi \emph{borī}, \textquotedblleft{}sack, coarse bag\textquotedblright{}, under it (CDIAL 9320).
\end{cousinweb}

\begin{grammarlens}[title={What you've built}]
One more word for this chapter.
\end{grammarlens}

\subsection*{Your turn}
\begin{itemize}
\raggedright
  \item \emph{Say it:} \emph{borī}
  \item \emph{Say it:} \emph{borī}, once more
  \item \emph{Say it:} \emph{ikk borī vich bahut kaṇak hundī hai.} — A sack holds a lot of wheat.
  \item \emph{Recall:} say the Punjabi for chopped straw, then the Punjabi for rice straw, then say \emph{borī} again
  \item \emph{Write it:} copy \textbf{\pa{ਬੋਰੀ}} once, checking the whole word
\end{itemize}

\begin{tcolorbox}[breakable,colback=teal!4,colframe=teal!35!black,title={Before you move on}]
What does \pa{ਬੋਰੀ} mean? (\textquotedblleft{}A sack\textquotedblright{}.) Say it once more.
\end{tcolorbox}
